\begin{opinion}

    En tesis anterior de la facultad se han propuesto algoritmos para encontrar el sistema compartimental de ecuaciones diferenciales lineal con respecto a los parámetros que mejor describa un conjunto de datos. Sin embargo, en todos estos trabajos se asumía conocida la dinámica poblacional detrás del fenómeno estudiado.

    El trabajo que culmina con esta tesis provee herramientas para determinar un grafo que refleje la dinámica poblacional del fenómeno descrito en los datos.
  
    Para realizar esta tesis Dario tuvo que incursionar en temas que no están incluidos en su plan de estudios; resolver creativamente problemas que aparecieron como parte del proceso de desarrollo  y  adquirir habilidades relacionadas con la escritura de documentos científicos y la presentación oral de los resultados.
  
    De este tiempo de trabajo quiero resaltar la seriedad, independencia con la que trabajó Darío, así como la calidad de las solucions propuestas por él.
  
     Creo que estamos en presencia de un trabajo excelente desarrollado por un excelente científico de la computación.
  
\begin{flushright}
    \underline{\hspace{6.5cm}}\\
    MSc. Fernando Raul Rodriguez Flores
 
    Facultad de Matemática y Computación
   
    Universidad de la Habana
 
    Junio, 2025
\end{flushright}

\end{opinion}